% !TEX root =  ../mainPPDP.tex
%!TEX spellcheck = en_GB


%In this section we introduce some notation to talk about communicating automata, executions and message sequence charts.

We assume a finite set of \emph{processes} $\Procs=\{p,q,\ldots\}$ and a finite set of messages (labels) 
$\Msg=\{\msg,\ldots\}$.
We consider two kinds of actions:  \emph{send actions} that are  of the form 
$\send{p}{q}{\msg}$ and are executed by  process $p$  sending message  $\msg$ to  $q$;
 \emph{receive actions} that are of the form $\recv{p}{q}{\msg}$ and are executed by  $q$ receiving  $\msg$ from  $p$.
We write $\Act$ for the finite set $\Procs\times\Procs\times\{!,?\}\times\Msg$ of all actions, and $\Actp$ for the set of actions 
that can be executed by $p$ (i.e., $\send{p}{q}{\msg}$ or $\recv{q}{p}{\msg}$). 
We omit processes when they are clear from the context and simply write $!\msg$ or $?\msg$ for a send or receive action, respectively.

An \emph{event} $\event$ of a sequence of actions $w\in \Act^*$, is an index $i$ in $\{1,\ldots,\length{w}\}$;
$i$ is a send (resp. receive) event of $w$ if $w[i]$ is a send (resp. receive) action.
We write $\sendeventsof{w}$ (resp. $\receiveeventsof{w}$) for the set of send (resp. receive) events 
of $w$ and $\eventsof{w}=\sendeventsof{w}\cup\receiveeventsof{w}$ for the set of all events of $w$.
When all events are labeled with distinct actions, we identify an event with its action.



\subsubsection*{Executions.}
An execution is a well defined sequence of actions $e\in\Act^*$, where 
a receive action is always preceded by a unique corresponding send action.


\begin{definition}[Execution]
	An \emph{execution} over $\Procs$ and $\Msg$ is a sequence of actions $e\in\Act^*$  
	where an injective function from receive events to send events $\source_{e}:\receiveeventsof{w}\to\sendeventsof{w}$ 
	can be defined such that for each receive event $i$ labeled with $\recv{p}{q}{\msg}$,
	 $\source_{e}(i)$ is labeled with $\send{p}{q}{\msg}$, and  $\source_{e}(i)<i$.
\end{definition}


%For two events $i,j\in\eventsof{e}$, we write $i\hbstrictof{e}{} j$ if $i<j$ (order on natural numbers).
%When all send actions of a sequence of actions $w$ are distinct, there is at most one execution 
%$\execution$ such that $w$ is the sequence of actions of $\execution$, and we often identify $w$ with $\execution$
%and let $\source$ be implicit. 

%An execution $\w_1$ is a \emph{prefix} of $w_2$ if $w_1$ is a prefix of $w_2$ and 
%$\source_1$ is the restriction of $\source_2$ to $\receiveeventsof{w_1}$.
For a set of executions $\mathcal{E}$, we write $\prefixclosureof{\mathcal{E}}$ 
for the set of all prefixes of the executions in $\mathcal{E}$. 
We say that an execution $e_2$ is a \emph{completion} of an execution $e_1$
if $e_1$ is a prefix of $e_2$. 
A \emph{concatenation} $e_1\cdot w$ of an execution
$e_1$ and a sequence of actions $w$ is the execution $e_2 = e_1\cdot w_2$ where 
$e_2$ is a completion of $e_1$ (note that $w$ is not an execution, since it may contain receive events
which sources are in $e_1$).  
The \emph{projection} $\projofon{e}{p}$ of an execution $e$ on a process $p$ 
is the subsequence of actions in $\Actp$.  
A send event $s$ is \emph{matched} if there is a receive event $r$ such that $s=\source(r)$.
An execution $e$ is \emph{orphan-free} if $\source$ is surjective over the send events of $e$, i.e.,
all send events are matched.

\subsubsection*{Communication Models.}
In this paper, we consider   two communication models: 
peer-to-peer ($\ppmodel$) and  synchronous ($\synchmodel$). 
However, as commented  in the conclusions,  this work is part of a bigger more general project and a large amount of the results of the paper can be extended to "any" communication model, or 
only require  few assumptions.  In this perspective, we  introduce a general definition of a communication model here
and discuss on specific results on how they can be extended to the general case.
\begin{definition}[Communication model]
	\label{def:communication-model}
	A \emph{communication model} $\acommunicationmodel$ is a set 
	$\executionsofmodel{\acommunicationmodel}$ of executions.
\end{definition}
% It is \emph{prefix-closed} if 
% $\executionsofmodel{\acommunicationmodel}$ is prefix-closed, i.e., for all
% $e,e'\in\executionsofmodel{\acommunicationmodel}$, if $e\prefixorder e'$, then $e\in\executionsofmodel{\acommunicationmodel}$.

In the  synchronous communication model $\synchmodel$, message exchanges can be thought as rendez-vous synchronisations.
In other words, an execution $e$ belongs to $\executionsofmodel{\synchmodel}$ if
all send events are immediately followed by their corresponding receive events.
\begin{definition}[Synchronous communication model]
	An execution $\execution=(w,\source) \in \executionsofmodel{\synchmodel}$ if
	for all send event $s\in\sendeventsof{e}$, $s+1$ is a receive event of $e$ and $\source(s+1)=s$.
\end{definition}

In the peer-to-peer communication model $\ppmodel$, messages sent by a process $p$ to $q$ transit over a FIFO channel
that is specific to the pair $(p,q)$: if $p$ sends first $m_1$ then $m_2$ to $q$, 
then $m_2$ cannot overtake $m_1$ in the FIFO channel. In particular:
\begin{itemize}
	\item if $m_1$ is not received, then $m_2$ is not received either;
	\item if both are received, then $m_1$ is received before $m_2$.
\end{itemize}
      
\begin{definition}[$\ppmodel$ communication model]\label{def:queue-based-communication-model}
	$\executionsofmodel{\ppmodel}$  is
	the set of executions $e$ such that for any two send events $s_1=\send{p}{q}{\msg_1}$ and 
	$s_2=\send{p}{q}{\msg_2}$ in $\sendeventsof{e}$, 
	with $s_1 < s_2$,
	one of the two holds:
	\begin{itemize}
		\item $s_2$ is unmatched, or
		\item there exists $r_1,r_2$ such that $r_1<r_2$, $\source(r_1)=s_1$, and $\source(r_2)=s_2$.
	\end{itemize}
\end{definition}

Note that, $\executionsofmodel{\synchmodel}\subset\executionsofmodel{\ppmodel}$.
Moreover, if $e$ is an execution in $\ppmodel$, then  $\source_e$ is defined as follows:
the source of the $i$-th receive event of $q$ from $p$ is the $i$-th send event of $p$ to $q$.
If $e$ is an execution in $\synchmodel$, then $\source_e$ is defined as follows: 
for all receive event $i$, $\source_{e}(i)=i-1$.


\subsubsection*{Message Sequence Charts.}
While executions correspond to a total order view of the events occurring in a 
system, message sequence charts (MSC) adopt a distributed, a partial order view on the events.
For a tuple $\msc=(w_p)_{p\in\Procs}$, each $w_p\in\Actsp$ 
is a sequence of actions representing  the ones executed by process $p$ according to some total, locally observable, order.
We write $\eventsof{\msc}$  for the set $\{(p,i) \mid p \in \Procs \text{ and } 0 \leq i < \length{w_p}\}$.
The label $\actionof{\event}$ of the event $\event=(p,i)$ is the action $w_p[i]$. The event $\event$
is a send (resp. receive) event if it is labeled with a send (resp. receive) action.
We write $\sendeventsof{\msc}$ (resp. $\receiveeventsof{\msc}$) for the set of send (resp. receive) events of $\msc$; we also write $\messageof{\event}$ for the message sent or received at event $\event$, and 
$\processof{\event}$ for the process executing $\event$. Finally, we write 
$\verticalorderstrict{\event_1}{\event_2}$ if there is a process $p$ and $i<j$ such that
$\event_1=(p,i)$ and $\event_2=(p,j)$.
%\pascal{Notations $\sendevents{\msc}$,  $\messageof{e}$, $\receiveevents{\msc}$ and $\processof{e}$ are used once in the paper. Should we keep them?}

\begin{definition}[MSC]\label{def:msc}
	An {MSC} over $\Procs$ and $\Msg$ is a tuple $\msc = \big((w_p)_{p\in\Procs},\source\big)$
    where
    \begin{enumerate}
        \item for each process $p$, $w_p\in\Actsp$ is a finite sequence of actions;
        \item $\source: \receiveeventsof{\msc} \to \sendeventsof{\msc}$ is 
            an injective function from receive events to send events such that
			for all receive event $\event$ labeled with $\recv{p}{q}{\msg}$,
            $\source(\event)$ is labeled with $\send{p}{q}{\msg}$.
    \end{enumerate}
\end{definition}

\begin{figure}[t]
		\captionsetup[subfigure]{justification=centering}
	% \centering
	\begin{subfigure}[t]{0.22\textwidth}\centering

		\begin{tikzpicture}[scale=0.7, every node/.style={transform shape}]
			\newproc{0}{p}{-2.2};
			\newproc{2}{q}{-2.2};

			\newmsgm{0}{2}{-1.7}{-0.5}{1}{0.25}{black};
			\newmsgm{2}{0}{-1.7}{-0.5}{2}{0.25}{black};

			\end{tikzpicture}
		\caption{non-linearisable MSC}	\label{fig:raw_msc_ex}

		\end{subfigure}
%        \begin{subfigure}[t]{0.24\textwidth}\centering
%
%            \begin{tikzpicture}[scale=0.7, every node/.style={transform shape}]
%                \newproc{0}{p}{-2.2};
%                \newproc{2}{q}{-2.2};
%    
%                \newmsgm{0}{2}{-0.5}{-1.5}{1}{0.1}{black};
%                \newmsgm{0}{2}{-1}{-1}{2}{0.8}{black};
%    
%                \end{tikzpicture}
%            \caption{$\bagmodel$, and $\erlangmodel$ if $m_1\neq m_2$}	\label{fig:bag_ex}
%    
%            \end{subfigure}
        % \centering
	\begin{subfigure}[t]{0.22\textwidth}\centering
		\begin{tikzpicture}[scale=0.7, every node/.style={transform shape}]
			\newproc{0}{p}{-2.2};
			\newproc{1}{q}{-2.2};
			\newproc{2}{r}{-2.2};

			\newmsgm{0}{1}{-0.3}{-2}{1}{0.15}{black};
			\newmsgm{0}{2}{-0.9}{-0.9}{2}{0.75}{black};
			\newmsgm{2}{1}{-1.5}{-1.5}{3}{0.5}{black};
			% \newmsgm{2}{1}{-2}{-2}{4}{0.5}{black};

			% \newflechevert{Purple}{0}{-0.3}{-0.9};
			% \newflechehor{Purple}{-0.9}{0}{2};
			% \newflechevert{Purple}{2}{-0.9}{-1.5};
		\end{tikzpicture}
		\caption{$\ppmodel$} \label{fig:pp_ex}
	\end{subfigure}
%	\begin{subfigure}[t]{0.24\textwidth}\centering
%		\begin{tikzpicture}[scale=0.7, every node/.style={transform shape}]
%			\newproc{0}{p}{-2.2};
%			\newproc{1}{q}{-2.2};
%			\newproc{2}{r}{-2.2};
%			\newproc{3}{s}{-2.2};
%	
%			\newmsgm{0}{3}{-0.6}{-0.6}{2}{0.8}{black};
%			\newmsgm{2}{3}{-1.1}{-1.1}{3}{0.5}{black};
%			\newmsgm{2}{1}{-1.5}{-1.5}{4}{0.5}{black};
%			\newmsgm{0}{1}{-0.3}{-2}{1}{0.5}{black};
%		\end{tikzpicture}
%		\caption{$\causalmodel$} \label{fig:causal_ex}
%	\end{subfigure}
	% \hfill
%\begin{subfigure}[t]{0.24\textwidth}\centering
%	\begin{tikzpicture}[scale=0.7, every node/.style={transform shape}]
%		\newproc{0}{p}{-2.2};
%		\newproc{1}{q}{-2.2};
%		\newproc{2}{r}{-2.2};
%		\newproc{3}{s}{-2.2};
%
%		\newmsgm{3}{0}{-0.6}{-0.6}{2}{0.2}{black};
%		\newmsgm{3}{2}{-1.1}{-1.1}{3}{0.5}{black};
%		\newmsgm{2}{1}{-1.5}{-1.5}{4}{0.5}{black};
%		\newmsgm{0}{1}{-0.3}{-2}{1}{0.5}{black};
%	\end{tikzpicture}
%			\caption{$\mbmodel$} \label{fig:mb_ex}
%\end{subfigure}
%\begin{subfigure}[t]{0.24\textwidth}\centering
%	\begin{tikzpicture}[scale=0.7, every node/.style={transform shape}]
%		\newproc{0}{p}{-2.2};
%		\newproc{1}{q}{-2.2};
%		\newproc{2}{r}{-2.2};
%
%		\newmsgm{0}{1}{-0.3}{-1.5}{1}{0.15}{black};
%		\newmsgm{1}{0}{-0.9}{-0.9}{2}{0.25}{black};
%		\newmsgm{1}{2}{-2}{-2}{3}{0.5}{black};
%		% \newmsgm{2}{1}{-2}{-2}{4}{0.5}{black};
%	\end{tikzpicture}
%		\caption{$\busmodel$} \label{fig:bus_ex}
%\end{subfigure}
\begin{subfigure}[t]{0.22\textwidth}\centering
	\begin{center}
		\begin{tikzpicture}[scale=0.7, every node/.style={transform shape}]
			\newproc{0}{p}{-2.2};
			\newproc{1}{q}{-2.2};
			\newproc{2}{r}{-2.2};
			\newproc{3}{s}{-2.2};

			\newmsgm{0}{1}{-0.5}{-0.5}{1}{0.5}{black};
			\newmsgm{2}{3}{-0.5}{-0.5}{2}{0.5}{black};
			\newmsgm{1}{2}{-1}{-1}{3}{0.5}{black};

		\end{tikzpicture}
		\caption{$\synchmodel$}
		\label{fig:rsc_ex}
	\end{center}
\end{subfigure}
\begin{subfigure}[t]{0.24\textwidth}\centering
	\begin{tikzpicture}[scale=0.7, every node/.style={transform shape}]
		\newproc{0}{p}{-2.2};
		\newproc{2}{q}{-2.2};

		\newmsgum{0}{2}{-0.8}{1}{0.2}{black};
		\newmsgm{0}{2}{-1.6}{-1.6}{2}{0.2}{black};

		\end{tikzpicture}
	\caption{MSC with an orphan message}	\label{fig:orphan_ex}
\end{subfigure}
		\caption{Examples of MSCs. The sender of a message is at the origin of the arrow and the receiver at the destination. Unmatched send events are depicted with dashed arrows.}\label{fig:exmscs}
\end{figure}

For an execution $\execution$,  $\mscof{\execution}$ is the MSC 
$\big((w_p)_{p\in\Procs},\source\big)$ where $w_p$ is the subsequence of $\execution$ restricted to the actions of $p$,
and $\source$ is the lifting of $\source_{\execution}$ to the events of $(w_p)_{p\in\Procs}$.


\begin{example}
    \label{ex:msc}
     MSC $\msc$ in Fig.~\ref{fig:raw_msc_ex} is an MSC over $\Procs=\{p,q\}$ and $\Msg=\{\msg_1,\msg_2\}$
    with $\msc=\large((w_p,w_q),\source\large)$, $w_p = ?\msg_2!\msg_1$, $w_q = ?\msg_1!\msg_2$, $\source((p,0)) = (q,1)$,
    and $\source((q,0)) = (p,1)$. Note that there is no execution $\execution$ such that 
	$\msc=\mscof{\execution}$ as all receptions precede the corresponding sends. On the other hand, the MSC
	in Fig.~\ref{fig:pp_ex} is $\mscof{\execution}$ for the execution $\execution$
	$=\large(!\msg_1!\msg_2?\msg_2!\msg_3?\msg_3?\msg_1,\source\large)$ where 
	$\source(3) = 2$, $\source(5) = 2$ and $\source(6) = 4$. It is the only execution that induces this MSC, but 
	in general there might exist several executions inducing the same MSC.
\end{example}

For a set of processes $\Procs$, 
an MSC $M=\big((w_p)_{p\in\Procs},\source\big)$ is the \emph{prefix} 
of another MSC $M'=\big((w'_p)_{p\in\Procs},\source'\big)$, in short $M \prefixorder M'$,
if for all $p\in\Procs$, $w_p$ is a prefix of $w'_p$ and $\source'(e)=\source(e)$ for all receive events $e$ of $M$.
The \emph{concatenation} of MSCs $M_1$ and $M_2$ is the MSC  $M_1\cdot M_2$ 
obtained by gluing "vertically"  $M_1$ before $M_2$:
formally, if $M_1=((w_p^1)_{p\in\Procs},\source_1)$ and $M_2=((w_p^2)_{p\in\Procs},\source_2)$, 
then $M_1\cdot M_2=((w_p)_{p\in\Procs},\source)$ where
%\begin{inparaenum}[(i)]
 %   \item 
    for all $p$, $w_p=w_p^1\cdot w_p^2$, and
 %   \item 
    $\source$ is defined by $\source(e) = \source_i(e)$ for all receive events $e$ of $M_i$ ($i=1,2$).
%\end{inparaenum}

\subsubsection*{Happens-before relation and linearisations}
%\etienne{TODO cette section et ailleurs: relire en faisant attention à la convention de notation entre ordres partiels $\prec$ et totaux $<$}
In a given MSC $M$,
an event $\event$ happens before  
another event $\event'$, if 
(i) $\event$ and $\event'$ are events
of a same process $p$ and happen in that order on 
the time line of $p$,
or (ii) $\event$ is send event matched by $\event'$,
or (iii) a sequence of such situations defines
a path from $\event$ to $\event'$.

\begin{definition}[Happens-before relation]
Let $M$ be an MSC. The happens-before relation over $M$
is the binary relation $\happensbeforestrict$ defined as 
the least transitive relation over $\eventsof{\msc}$ such that:
\begin{itemize}
   \item 
   for all 
   $p,i,j$, if $i<j$, then $(p,i)\happensbeforestrict (p,j)$, and
   \item 
   for all receive events $\event$, $\source(\event) \happensbeforestrict \event$.
\end{itemize}
\end{definition}

\begin{definition}[Linearisation]
	\label{def:linearisation}
	A \emph{linearisation} of an MSC $\msc$ is a
total order $\alinearisation$ on $\eventsof{\msc}$
	that refines $\happensbeforestrict$:  for all events $\event,\event'$, 
	if $\event\happensbeforestrict \event'$, then $\event\alinearisation \event'$. 
\end{definition}
We write $\linearisationsof{\msc}{}$ for the set of all linearisations of $\msc$.
We often identify a linearisation
with the execution it induces. 

\begin{example}
	\label{ex:linearisation}
	Let $\msc$ be the MSC in Fig.~\ref{fig:rsc_ex}. 
	Then $!m_1$ happens before $?m_1$, which
	happens before $!m_3$, and both $!m_3$ and $!m_2$
	happen before $?m_3$.
	Moreover, $\happensbeforestrict$ is a partial order, 
	and 
	$!m_1!m_2?m_1!m_3?m_3?m_2 \in \linearisationsof{\msc_c}{}$.
	On the other hand, consider the MSC $M$ in 
	in Fig.~\ref{fig:raw_msc_ex}; then
	$\happensbeforestrictof{\msc'}$ is not a partial order, because
	$?m_2\happensbeforestrictof{\msc'}!m_1\happensbeforestrictof{\msc'}?m_1\happensbeforestrictof{\msc'}!m_2\happensbeforestrictof{\msc'}?m_2$,
	therefore $\linearisationsof{\msc'}{}=\emptyset$.
\end{example}


Given an MSC $\msc$, we write 
$\linearisationsof{\msc}{\acommunicationmodel}$
to denote 
$\linearisationsof{\msc}{}\cap\executionsofmodel{\acommunicationmodel}$; 
the executions of $\linearisationsof{\msc}{\acommunicationmodel}$
are called the linearisations of $\msc$ 
in the communication model $\acommunicationmodel$. 	

\begin{definition}[$\acommunicationmodel$-linearisable MSC]
	\label{def:linearisable-msc}
	An MSC $\msc$ is \emph{linearisable} in a communication model $\acommunicationmodel$ if 
	$\linearisationsof{\msc}{\acommunicationmodel}\neq\emptyset$.
	We write $\mscsetofmodel{\acommunicationmodel}$ for the set of all MSCs linearisable in $\acommunicationmodel$.
\end{definition}

\begin{example}
	\label{ex:msc-linearisable}
	The MSC $M_b$ in Fig.~\ref{fig:pp_ex} is linearisable in $\ppmodel$,
	with $\linearisationsof{M_b}{\ppmodel}=
%	$ 
%	$$
		\{!m_1!m_2?m_2!m_3?m_3?m_1\}
		$.
%	$$
	However, $M_b$ is not linearisable in $\synchmodel$.
	Finally the MSC $M_c$ in Fig.~\ref{fig:rsc_ex} is linearisable in $\synchmodel$ with
	$\linearisationsof{M_c}{\synchmodel}=$
	$$
	\{~!m_1?m_1!m_2?m_2!m_3?m_3~,~
	   !m_2?m_2!m_1?m_1!m_3?m_3~\}
	$$
\end{example}

Finally, we introduce a property that will be helpful in the next paragraph 
for
giving an alternative characterisation of deadlock-freedom of a system of communicating finite state machines.

\begin{definition}[Causally-closed communication model]\label{def:causally-closed-communication-model}
	A communication model $\acommunicationmodel$ is \emph{causally-closed} if for all MSCs $M$,
	$\linearisationsof{\msc}{\acommunicationmodel}\neq\emptyset$ implies that
	$\linearisationsof{\msc}{\acommunicationmodel}=\linearisationsof{\msc}{}$.
\end{definition}


  
  % !TEX root =  ../mainPPDP.tex
%!TEX spellcheck = en_GB
Observe that not all communication models are causally closed. It is the case for $\ppmodel$,  but it is immediate  to see that the property is not valid  for $\synchmodel$. Take for instance  MSC $M$ in Fig.~\ref{fig:rsc_ex},  its linearisation 
$!m_1!m_2?m_1?m_2!m_3?m_3$ 
  does not belong to $\linearisationsof{M_c}{\synchmodel}$.  To show that $\ppmodel$ is causally closed, 
notice that $<$ can be replaced with $\porderof{M}$ in  Definition~\ref{def:queue-based-communication-model}.
\begin{lemma}\label{lem:reformulation-of-ppmodel-def}
    Let $e$ be an execution and  $M=\mscof{e}$.
    Then $e$ is an execution of $\ppmodel$ if and only if
    for any two send events $s_1=\send{p}{q}{\msg_1}$ and $s_2=\send{p}{q}{\msg_2}$ in $\sendeventsof{e}$, with $s_1 \porderof{M} s_2$,
        one of the two holds:
        \begin{itemize}
            \item $s_2$ is unmatched, or
            \item there exists $r_1,r_2$ such that $r_1 \porderof{M} r_2$, $\source(r_1)=s_1$, and $\source(r_2)=s_2$.
        \end{itemize}
\end{lemma}
\begin{proof}
    Assume that $e$ is a $\ppmodel$ execution.
    Let $s_1=\send{p}{q}{\msg_1}$ and $s_2=\send{p}{q}{\msg_2}$ be two send events in $\sendeventsof{e}$ such that $s_1 \porderof{M} s_2$.
    Then $s_1<s_2$ in $e$, because $<$ is a linearisation of $\porderof{M}$.
    By  definition of $\ppmodel$, $s_2$ is unmatched or there exists $r_1,r_2$ such that $r_1 < r_2$, $\source(r_1)=s_1$, and $\source(r_2)=s_2$.
    In the second case, $r_1 < r_2$ implies that
    $r_1 \porderof{M} r_2$, because both $r_1$ and $r_2$ occur on the same process $q$.
    Conversely, assume that $e$ and $M$ satisfy the above reformulation of the definition of $\ppmodel$. 
    Let $s_1=\send{p}{q}{\msg_1}$ and $s_2=\send{p}{q}{\msg_2}$ be two send events in $\sendeventsof{e}$ such that $s_1 < s_2$.
    Then $s_1 \porderof{M} s_2$
    because $s_1$ and $s_2$ occur on the same process $p$.
    By the reformulation of the definition of $\ppmodel$, $s_2$ is unmatched or there exists $r_1,r_2$ such that $r_1 \porderof{M} r_2$, $\source(r_1)=s_1$, and $\source(r_2)=s_2$.
    In the second case, $r_1 \porderof{M} r_2$
    implies that $r_1 < r_2$ because $<$ is a  linearisation of $\porderof{M}$.
\end{proof}

\begin{restatable}[]{lemma}{lemppiscasuallyclosed}\label{lem:pp-is-causally-closed}
	$\executionsofmodel{\ppmodel}$ is causally-closed.
\end{restatable}
\begin{proof}
    Let $M\in\mscsetofmodel{\ppmodel}$ 
    and  $e$ be a linearisation of $M$.
    We want to show that $e\in\executionsofmodel{\ppmodel}$.
    By definition of $\mscsetofmodel{\ppmodel}$, 
    there is an execution 
    $e'\in\executionsofmodel{\ppmodel}$
    such that $\mscof{e'}=M$.
    By Lemma~\ref{lem:reformulation-of-ppmodel-def},
    $e$ is also a $\ppmodel$ execution.
\end{proof}
  %The proof of the following Lemma is in Appendix \ref{app:pp-is-causally-closed}.




\subsubsection*{Communicating finite state machines.}
We assume standard notations for automata, words and languages. As usual, a non-deterministic finite state automaton (NFA) is a tuple
$\mathcal A = (Q,\Sigma, \delta, l_0, F)$ where $Q$ is a set of control states, 
$\Sigma$ is an alphabet, $\delta:Q\times\Sigma\to 2^Q$ 
is the transition relation,
$l_0$ is the initial control state, and $F\subseteq Q$ is the set of accepting states.
The language $\languageofnfa{\mathcal A}$ of an NFA $\mathcal A$ and the notion of deterministic
finite state automaton (DFA) or $\varepsilon$ transitions are defined as usual.
We write $\acceptcompletion{\mathcal A}{}$ for the automaton obtained from $\mathcal A$ by setting
$F=Q$. We start  by recalling the definition of communicating finite state machines~\cite{BrandZafiropulo}.

\begin{definition}[CFSM] 
    A communicating finite state machine (CFSM) is an NFA with $\varepsilon$-transitions $\acfsm$ over the alphabet $\Act$.
    A system of CFSMs is a tuple $\cfsms = (\acfsm_p)_{p\in\Procs}$.
\end{definition}

 Given a system of CFSMs $\cfsms=(\acfsm_p)_{p\in\Procs}$,
we write $\acceptcompletion{\cfsms}{}$ for the system of CFSMs $\acceptcompletion{\cfsms}=(\acceptcompletion{\acfsm_p})_{p\in\Procs}$
where all states are accepting, i.e., $F_p = Q_p$.

\begin{definition}[Executions of a CFSMs in $\acommunicationmodel$]\label{def:executions-of-cfsms}
Given a system 
$\cfsms = (\acfsm_p)_{p\in\Procs}$ of CFSMs, and a communication model $\acommunicationmodel$,
 $\executionsofcfsms{\cfsms}{\acommunicationmodel}$ is the set  of \emph{executions} $e\in\executionsofmodel{\acommunicationmodel}$
such that for all $p$, $\projofon{e}{p}$ is in $\languageofnfa{\acfsm_p}{}$.
\end{definition}

\begin{remark}
	Let $\acommunicationmodel$ be a communication model, 
	$\cfsms$  a system of CFSMs, and $e,e'\in\executionsofmodel{\acommunicationmodel}$ such that $\mscof{e}=\mscof{e'}$, 
	then $e\in\executionsof{\cfsms}{\acommunicationmodel}$ 
	if and only if $e'\in\executionsof{\cfsms}{\acommunicationmodel}$. This follows from the fact that $\projofon{e}{p}= \projofon{e'}{p}$ for all $p$.
\end{remark}

We write $\mscsofcfsms{\cfsms}{\acommunicationmodel}$ for the set $\{\mscof{e} \mid e\in\executionsofcfsms{\cfsms}{\acommunicationmodel}\}$.

A system is orphan-free if, whenever all machines have reached an accepting state, no message
remains in transit, i.e., no message is sent but not received.

\begin{definition}[Orphan-free]\label{def:orphan-free}
	A system of CFSMs $\cfsms$ is \emph{orphan-free} for a communication model 
	$\acommunicationmodel$ if
	for all $e\in\executionsofcfsms{\cfsms}{\acommunicationmodel}$,
	$e$ is orphan-free.
\end{definition}

All  synchronous executions are orphan-free by definition.

A system is deadlock-free if, 
any \emph{partial} execution can be extended/completed to an accepting execution.

\begin{definition}[Deadlock-free]\label{def:deadlock-free}
	A system of CFSMS $\cfsms$ is \emph{deadlock-free} for a communication model
	$\acommunicationmodel$ if for all 
	$e\in\executionsofcfsms{\acceptcompletion{\cfsms}}{\acommunicationmodel}$,
	there is an execution $e'\in\executionsofcfsms{\cfsms}{\acommunicationmodel}$ 
	such that $e\prefixorder e'$.
\end{definition}

\begin{remark}\label{rem:equivalent-formulation-of-deadlock-free-cfsms}
	A system of CFSMS $\cfsms$ is \emph{deadlock-free} for a communication model
	$\acommunicationmodel$ if and only if 
	$$
	\executionsofcfsms{\acceptcompletion{\cfsms}}{\acommunicationmodel}\subseteq
	\prefixclosureof{\executionsofcfsms{\cfsms}{\acommunicationmodel}}
	$$
\end{remark}

The following result shows that, for either $\ppmodel$ or $\synchmodel$ communication
models, the deadlock-freedom property of a system of CFSMs can be expressed as a property on the MSCs
of the system. %The proof can be found in Appendix \ref{app:deadlock-free-as-a-property-on-mscs-for-p2p-and-synch}.

\begin{restatable}[Deadlock-freedom as a MSC property]{proposition}{propdeadlockfreeasapropertyonmscsforppandsynch}
	\label{prop:deadlock-free-as-a-property-on-mscs-for-p2p-and-synch}
	Assume $\acommunicationmodel$ is $\ppmodel$ (respectively, $\acommunicationmodel$ is $\synchmodel$).
	Then a system of CFSMs $\cfsms$ is deadlock-free for $\acommunicationmodel$ if and only if
	$\mscsofcfsms{\acceptcompletion\cfsms}{\acommunicationmodel}\subseteq\prefixclosureof{\mscsofcfsms{\cfsms}{\acommunicationmodel}}$.
\end{restatable}
% !TEX root =  ../mainPPDP.tex
%!TEX spellcheck = en_GB

\begin{proof}
    We show each implication separately.
    \begin{description}
        \item[$\Rightarrow$:]
        if $\cfsms$ is deadlock-free: $\executionsof{\acceptcompletion{\cfsms}}{\acommunicationmodel}\subseteq\prefixclosureof{\executionsof{\cfsms}{\acommunicationmodel}}$, 
        then 
        $\msclanguageof{\acceptcompletion{\cfsms}}{\acommunicationmodel}\subseteq
        \prefixclosureof{\msclanguageof{\cfsms}{\acommunicationmodel}}$.
        For this implication, $\acommunicationmodel$ can be any communication model.
        Let $\msc\in\msclanguageof{\acceptcompletion{\cfsms}}{\acommunicationmodel}$, 
       we show that $\msc\in\prefixclosureof{\msclanguageof{\cfsms}{\acommunicationmodel}}$.
        By definition, there is an execution $e\in\executionsof{\acceptcompletion{\cfsms}}{\acommunicationmodel}$ such that $\msc=\mscof{e}$.
        By hypothesis, there is a completion $e' \in\executionsof{\cfsms}{\acommunicationmodel}$ of $e$.
        By definition, $\mscof{e'}\in\msclanguageof{\acceptcompletion{\cfsms}}{\acommunicationmodel}$,
        so $\msc\in\prefixclosureof{\msclanguageof{\cfsms}{\acommunicationmodel}}$.

        \item[$\Leftarrow$:] for  $\acommunicationmodel=\ppmodel$:
        if  
        $\msclanguageof{\acceptcompletion{\cfsms}}{\acommunicationmodel}\subseteq
        \prefixclosureof{\msclanguageof{\cfsms}{\acommunicationmodel}}$,
        then $\cfsms$ is deadlock-free.
        Let $e\in\executionsof{\acceptcompletion{\cfsms}}{\acommunicationmodel}$ be an execution,
        wez show that $e$ has a completion in $\executionsof{\cfsms}{\acommunicationmodel}$.
        By definition, $\mscof{e}\in\msclanguageof{\acceptcompletion{\cfsms}}{\acommunicationmodel}$.
        By hypothesis, there is a MSC $\msc'\in\msclanguageof{\cfsms}{\acommunicationmodel}$
        such that $\mscof{e}\prefixordermsc\msc'$. By definition of $\prefixordermsc$ on MSCs,
        there are two executions $e_1,e'$ such that 
        $e_1\prefixorder e'$, $\mscof{e'}=\msc'$, 
        and $\mscof{e_1}=\mscof{e}$.
        Let $<$ be the binary relation on $\eventsof{M'}$ 
        defined as $< \eqdef \torderstrictof{e}\cup <_1\cup <_2$, where:
        \begin{itemize}
            \item $\event <_1<\event'$ if $\event\not\in\eventsof{M}$,
            $\event'\not\in\eventsof{M}$ and $\event <_{e'} \event'$;
            \item $\event <_2\event'$ if $\event\in\eventsof{M}$ and $\event'\not\in\eventsof{M}$.
        \end{itemize}
        Note that $<_1$ and $<_2$ are transitive. We claim that $<$
        is also transitive. Indeed, 
        $$
        \begin{array}{rcl}
           \torderof{e}\cdot <_1 & = & <_1\\
            <_1\cdot <_2 & = & <_1 \\
            \torderof{e}\cdot<_2 & = & \emptyset \\
            (<_1\cup <_2)\cdot \torderof{e} & = & \emptyset \\
            (\torderof{e}\cup <_2)\cdot <_1 & = & \emptyset \\
        \end{array}
        $$
        Moreover, $<$ is irreflexive, because $\torderof{e}$ is irreflexive and $<_1$ and $<_2$ are irreflexive.
        Let $\leq$ denote the reflexive closure of $<$.
        Then this is a partial order on $\eventsof{M'}$.
        By construction, $\leq$ is actually a total order on $\eventsof{M'}$.       
        Finally, we claim that $\leq$ is a linearisation of $M'$, i.e., $\porderof{M'}\subseteq \leq$.
        Indeed, assume that $\event\porderstrictof{M'}\event'$, we show that $\event< \event'$:
        \begin{itemize}
            \item if $\event\in\eventsof{M}$ and $\event'\in\eventsof{M}$,
            then $\event\porderstrictof{M} \event'$ (as $M$ is a prefix of $M'$), hence $\event\torderof{e} \event'$ 
            (because $e$ is a linearisation of $M$),
            and finally $\event < \event'$ by definition of $<$.
            
            \item if $\event\not\in\eventsof{M}$ and $\event'\not\in\eventsof{M}$,
            then $\event\torderof{e'} \event'$ (because $e'$ is a linearisation of $M'$), 
            therefore $\event <_1 \event'$ by definition of $<_1$, and finally $\event < \event'$ by definition of $<$.

            \item if $\event\in\eventsof{M}$ and $\event'\not\in\eventsof{M'}$,
            then $\event<_2 \event'$ by definition of $<_2$, therefore $\event < \event'$ by definition of $<$.
        \end{itemize}
        Therefore, $\leq$ is a linearisation of $M'$. Let $e''$ be the execution associated to $\leq$.
        Then
        \begin{itemize}
            \item $e\prefixorder e''$ by definition of $\leq$;
            \item $\mscof{e''}=M'$ as $\leq$ is a linearisation of $M'$;
            \item $e''\in\executionsof{\cfsms}{\acommunicationmodel}$ because $e''$ is a linearisation of $M'$, $M'\in\msclanguageof{\cfsms}{\acommunicationmodel}$,
            and $\ppmodel$ is causally closed (Lemma~\ref{lem:pp-is-causally-closed}).
        \end{itemize}
        Altogether, we have shown that $e$ has a completion in $\executionsof{\cfsms}{\acommunicationmodel}$, which ends the proof of 
        the reverse implication.

        \item[$\Leftarrow$:]   for $\acommunicationmodel=\synchmodel$:
            the proof is similar to previous case. We define the linearisation $\leq$ of $M'$ 
            in the exactly  same way, but we cannot use the fact that $\acommunicationmodel$ is causally closed.
            Instead, we use the fact that both $<_e$ and  $<_1$ are synchronous linearisations,
            and therefore $<$ is a synchronous linearisation as the concatenation of two synchronous linearisations.

    \end{description}
\end{proof}

%\iflong
%% !TEX root =  ../mainPPDP.tex
%!TEX spellcheck = en_GB

\begin{proof}
    We show each implication separately.
    \begin{description}
        \item[$\Rightarrow$:]
        if $\cfsms$ is deadlock-free: $\executionsof{\acceptcompletion{\cfsms}}{\acommunicationmodel}\subseteq\prefixclosureof{\executionsof{\cfsms}{\acommunicationmodel}}$, 
        then 
        $\msclanguageof{\acceptcompletion{\cfsms}}{\acommunicationmodel}\subseteq
        \prefixclosureof{\msclanguageof{\cfsms}{\acommunicationmodel}}$.
        For this implication, $\acommunicationmodel$ can be any communication model.
        Let $\msc\in\msclanguageof{\acceptcompletion{\cfsms}}{\acommunicationmodel}$, 
       we show that $\msc\in\prefixclosureof{\msclanguageof{\cfsms}{\acommunicationmodel}}$.
        By definition, there is an execution $e\in\executionsof{\acceptcompletion{\cfsms}}{\acommunicationmodel}$ such that $\msc=\mscof{e}$.
        By hypothesis, there is a completion $e' \in\executionsof{\cfsms}{\acommunicationmodel}$ of $e$.
        By definition, $\mscof{e'}\in\msclanguageof{\acceptcompletion{\cfsms}}{\acommunicationmodel}$,
        so $\msc\in\prefixclosureof{\msclanguageof{\cfsms}{\acommunicationmodel}}$.

        \item[$\Leftarrow$:] for  $\acommunicationmodel=\ppmodel$:
        if  
        $\msclanguageof{\acceptcompletion{\cfsms}}{\acommunicationmodel}\subseteq
        \prefixclosureof{\msclanguageof{\cfsms}{\acommunicationmodel}}$,
        then $\cfsms$ is deadlock-free.
        Let $e\in\executionsof{\acceptcompletion{\cfsms}}{\acommunicationmodel}$ be an execution,
        wez show that $e$ has a completion in $\executionsof{\cfsms}{\acommunicationmodel}$.
        By definition, $\mscof{e}\in\msclanguageof{\acceptcompletion{\cfsms}}{\acommunicationmodel}$.
        By hypothesis, there is a MSC $\msc'\in\msclanguageof{\cfsms}{\acommunicationmodel}$
        such that $\mscof{e}\prefixordermsc\msc'$. By definition of $\prefixordermsc$ on MSCs,
        there are two executions $e_1,e'$ such that 
        $e_1\prefixorder e'$, $\mscof{e'}=\msc'$, 
        and $\mscof{e_1}=\mscof{e}$.
        Let $<$ be the binary relation on $\eventsof{M'}$ 
        defined as $< \eqdef \torderstrictof{e}\cup <_1\cup <_2$, where:
        \begin{itemize}
            \item $\event <_1<\event'$ if $\event\not\in\eventsof{M}$,
            $\event'\not\in\eventsof{M}$ and $\event <_{e'} \event'$;
            \item $\event <_2\event'$ if $\event\in\eventsof{M}$ and $\event'\not\in\eventsof{M}$.
        \end{itemize}
        Note that $<_1$ and $<_2$ are transitive. We claim that $<$
        is also transitive. Indeed, 
        $$
        \begin{array}{rcl}
           \torderof{e}\cdot <_1 & = & <_1\\
            <_1\cdot <_2 & = & <_1 \\
            \torderof{e}\cdot<_2 & = & \emptyset \\
            (<_1\cup <_2)\cdot \torderof{e} & = & \emptyset \\
            (\torderof{e}\cup <_2)\cdot <_1 & = & \emptyset \\
        \end{array}
        $$
        Moreover, $<$ is irreflexive, because $\torderof{e}$ is irreflexive and $<_1$ and $<_2$ are irreflexive.
        Let $\leq$ denote the reflexive closure of $<$.
        Then this is a partial order on $\eventsof{M'}$.
        By construction, $\leq$ is actually a total order on $\eventsof{M'}$.       
        Finally, we claim that $\leq$ is a linearisation of $M'$, i.e., $\porderof{M'}\subseteq \leq$.
        Indeed, assume that $\event\porderstrictof{M'}\event'$, we show that $\event< \event'$:
        \begin{itemize}
            \item if $\event\in\eventsof{M}$ and $\event'\in\eventsof{M}$,
            then $\event\porderstrictof{M} \event'$ (as $M$ is a prefix of $M'$), hence $\event\torderof{e} \event'$ 
            (because $e$ is a linearisation of $M$),
            and finally $\event < \event'$ by definition of $<$.
            
            \item if $\event\not\in\eventsof{M}$ and $\event'\not\in\eventsof{M}$,
            then $\event\torderof{e'} \event'$ (because $e'$ is a linearisation of $M'$), 
            therefore $\event <_1 \event'$ by definition of $<_1$, and finally $\event < \event'$ by definition of $<$.

            \item if $\event\in\eventsof{M}$ and $\event'\not\in\eventsof{M'}$,
            then $\event<_2 \event'$ by definition of $<_2$, therefore $\event < \event'$ by definition of $<$.
        \end{itemize}
        Therefore, $\leq$ is a linearisation of $M'$. Let $e''$ be the execution associated to $\leq$.
        Then
        \begin{itemize}
            \item $e\prefixorder e''$ by definition of $\leq$;
            \item $\mscof{e''}=M'$ as $\leq$ is a linearisation of $M'$;
            \item $e''\in\executionsof{\cfsms}{\acommunicationmodel}$ because $e''$ is a linearisation of $M'$, $M'\in\msclanguageof{\cfsms}{\acommunicationmodel}$,
            and $\ppmodel$ is causally closed (Lemma~\ref{lem:pp-is-causally-closed}).
        \end{itemize}
        Altogether, we have shown that $e$ has a completion in $\executionsof{\cfsms}{\acommunicationmodel}$, which ends the proof of 
        the reverse implication.

        \item[$\Leftarrow$:]   for $\acommunicationmodel=\synchmodel$:
            the proof is similar to previous case. We define the linearisation $\leq$ of $M'$ 
            in the exactly  same way, but we cannot use the fact that $\acommunicationmodel$ is causally closed.
            Instead, we use the fact that both $<_e$ and  $<_1$ are synchronous linearisations,
            and therefore $<$ is a synchronous linearisation as the concatenation of two synchronous linearisations.

    \end{description}
\end{proof}
%\else
%The proof of Proposition~\ref{prop:deadlock-free-as-a-property-on-mscs-for-p2p-and-synch} can be found in Appendix~\ref{app:deadlock-free-as-a-property-on-mscs-for-p2p-and-synch}.
%\fi
